% الجزء: sets-functions-relations
% الفصل: relations
% القسم: orders

\documentclass[../../../include/open-logic-section]{subfiles}

\begin{document}

\olfileid[ar]{sfr}{rel}{ord}
\olsection{الترتيبات}

\begin{explain}
كثيرًا ما نقارن الأشياء فنقول إن هذا «أصغر من» ذاك أو «مساو له»
أو «أكبر منه» من جهة ما؛ وتلك علاقات \emph{ترتيب}. وليس الترتيب
نوعًا واحدًا: فمنه ما يوجب إمكان المقارنة بين كل شيئين، ومنه ما لا
يوجبه؛ ومنه ما يدخل الهوية، كـ~$\le$، وما يخرجها، كـ~$<$.
فلنميز هذه الأنواع بتصنيفها.
\end{explain}

\begin{defn}[الترتيب المسبق]
\emph{الترتيب المسبق} العلاقة الجامعة للانعكاسية والتعدية.
\end{defn}

\begin{defn}[الترتيب الجزئي]
\emph{الترتيب الجزئي} ترتيب مسبق يضاف إلى خصائصه مضادّة التناظر.
\end{defn}

\begin{defn}[الترتيب الخطي]\ollabel{def:linearorder}
\emph{الترتيب الكلي}، ويسمى أيضًا \emph{الترتيب الخطي}، ترتيب جزئي متصل.
\end{defn}

\begin{ex}
فالخطي جزئي، والجزئي مسبق، ولا ينعكس واحد من القولين.
فالعلاقة الشاملة على~$A$ انعكاسية ومتعدية، فهي ترتيب مسبق.
فإن كان لـ$A$ أكثر من !!{element} واحد انتفت مضادّة التناظر،
فلم تكن ترتيبًا جزئيًا.
\end{ex}

\begin{ex}
ولننظر في علاقة \emph{لا يزيد طولًا على} $\preccurlyeq$ على~$\Bin^*$:
$x \preccurlyeq y$ إذا وفقط إذا كان $\len{x} \le \len{y}$.
فهي ترتيب مسبق لانعكاسيتها وتعديتها، وهي متصلة كذلك. ولكنها ليست
جزئية، إذ لا تضاد التناظر: فـ$01 \preccurlyeq 10$ و$10 \preccurlyeq 01$
مع $01 \neq 10$.
\end{ex}

\begin{ex}
ومن أهم الترتيبات الجزئية $\subseteq$ على مجموعة من المجموعات.
ولا يلزم أن تكون خطية: فإذا كان $a \neq b$، ففي
$\Pow{\{a, b\}} = \{\emptyset, \{a\}, \{b\}, \{a,b\}\}$
نجد $\{a\} \nsubseteq \{b\}$ و$\{a\} \neq \{b\}$ و$\{b\}
\nsubseteq \{a\}$.
\end{ex}

\begin{ex}
ومن أمثلتها أيضًا \emph{القابلية للقسمة دون باق}. فإذا كان $n$ و~$m$
صحيحين، فمعنى $n \mid m$ أن $n$ يقسم $m$ بلا باق؛ أي يوجد صحيح~$k$
يحقق $m = kn$. وهذه على~$\Nat$ ترتيب جزئي غير خطي، إذ
$2 \nmid 3$ و$3 \nmid 2$. وأما على $\Int$ فليست إلا ترتيبًا مسبقًا:
يمتنع وصفها بمضادّة التناظر، لأن $1 \mid -1$ و$-1 \mid 1$ مع
$1 \neq -1$.
\end{ex}

\begin{ex}
وعلاقة \emph{الامتداد} على~$A^*$ تعريفها: $s \sqsubseteq s'$ إذا وفقط
إذا كان $s = \emptyseq$، أي المتتالية الخالية، أو $s = s'$، أو كان
$s = \tuple{s_1, \dots, s_n}$ و$s' = \tuple{s_1, \dots, s_n, s_{n+1}, \dots, s_m}$.
فمتى كان $s \sqsubseteq s'$ سمينا $s$ \emph{مقطعًا ابتدائيًا} من~$s'$.
وهي على $A^*$ ترتيب جزئي، لا خطي عمومًا؛ فإذا كان $a \neq b$ كان
$ab \not\sqsubseteq ba$ و$ba \not\sqsubseteq ab$.
\end{ex}

\begin{defn}[الترتيب الصارم]
كل علاقة تجمع اللاانعكاسية واللاتناظرية والتعدية \emph{ترتيب صارم}.
\end{defn}

\begin{defn}[الترتيب الخطي الصارم]\ollabel{def:strictlinearorder}
\emph{الترتيب الكلي الصارم}، أو \emph{الخطي الصارم}، ترتيب صارم متصل.
\end{defn}

\begin{ex}
يقابل $\le$، وهو خطي، الترتيب الخطي الصارم~$<$.
ويقابل $\subseteq$، وهو جزئي، الترتيب الصارم~$\subsetneq$.
\end{ex}

إذا أخذنا ترتيبًا صارمًا $R$ على~$A$ ثم ضممنا إليه القطر $\Id{A}$،
أي جميع الأزواج~$\tuple{x, x}$، صار ترتيبًا جزئيًا. وهذا
\emph{الإغلاق الانعكاسي} لـ~$R$. وبالعكس، حذف~$\Id{A}$ من ترتيب
جزئي يعطي ترتيبًا صارمًا. وتفصيل القول في القضيتين الآتيتين.

\begin{prop}\ollabel{prop:stricttopartial}
ليكن $R$ ترتيبًا صارمًا على~$A$. فإن $R^+ = R \cup \Id{A}$ ترتيب
جزئي؛ وإن كان $R$ خطيًا صارمًا كان $R^+$ خطيًا.
\end{prop}

\begin{proof}
لنفرض $R$ ترتيبًا صارمًا؛ أي $R \subseteq A^2$، و$R$ لا انعكاسية
ولا تناظرية ومتعدية. ضع $R^+ = R \cup \Id{A}$. فالمطلوب إثبات
انعكاسية $R^+$ ومضادّتها للتناظر وتعديتها.

أما انعكاسية $R^+$ فظاهرة من $\tuple{x, x} \in \Id{A} \subseteq
R^+$ لكل $x \in A$.

وأما مضادّة التناظر لـ$R^+$، فلو كان $R^+xy$ و$R^+yx$ مع
$x \neq y$، لكان $\tuple{x,y} \in R \cup \Id{A}$، ولما كان
$\tuple{x, y} \notin \Id{A}$ لزم $\tuple{x, y} \in R$، أي
$Rxy$. وبمثله~$Ryx$. وهذا يناقض لاتناظرية $R$ المفروضة.

وأما التعدية، فافرض $R^+xy$ و$R^+yz$. فإن كان
$\tuple{x, y} \in R$ و$\tuple{y,z} \in R$ لزم $\tuple{x, z} \in
R$ من تعدية~$R$. وإن لم يجتمعا، فلا بد من
$\tuple{x, y} \in \Id{A}$، أي $x = y$، أو $\tuple{y, z} \in \Id{A}$،
أي $y = z$. ففي الأولى $R^+yz$ بالفرض، ومع $x = y$ ينتج
$R^+xz$. والثانية مثلها. ففي كل حال $R^+xz$، فتثبت تعدية $R^+$.

ويبقى حكم الخطية. إن كانت $R$ متصلة، كان لكل $x \neq y$ أحد
$Rxy$ و~$Ryx$، أي أحد $\tuple{x, y} \in R$ و$\tuple{y, x} \in R$.
ولأن $R \subseteq R^+$ يبقى ذلك في $R^+$، فتكون $R^+$ متصلة أيضًا.
\end{proof}

\begin{prop}\ollabel{prop:partialtostrict}
ليكن $R$ ترتيبًا جزئيًا على~$A$. فإن $R^- = R \setminus \Id{A}$
ترتيب صارم؛ وإن كان $R$ خطيًا كان $R^-$ خطيًا صارمًا.
\end{prop}

\begin{proof}
إثبات هذه القضية موكول إلى التمرين.
\end{proof}

\begin{prob}
برهن على \olref[sfr][rel][ord]{prop:partialtostrict}.
\end{prob}

ومن خصائص الترتيب الخطي الصارم ما يشبه الامتدادية، وتبينه هذه النتيجة اليسيرة:

\begin{prop}\ollabel{prop:extensionality-strictlinearorders}
لكل ترتيب خطي صارم $<$ على $A$ تصدق الصيغة:
\[
  (\forall a, b \in A)((\forall x \in A)(x < a \liff x < b) \lif a = b).
\]
\end{prop}

\begin{proof}
لنفرض $(\forall x \in A)(x < a \liff x < b)$. فلو كان $a < b$ للزم
$a < a$، خلافًا للاانعكاسية $<$؛ إذن $a \nless b$.
وبمثله $b \nless a$. فلا يبقى إلا $a = b$، لاتصالية $<$.
\end{proof}

\end{document}
